```latex
\begin{theorem}
The number $\sqrt{2}$ is irrational.
\end{theorem}

\begin{proof}
Suppose, for the sake of contradiction, that $\sqrt{2}$ is rational. Then there exist integers $p$ and $q$ with $q \neq 0$ such that
\[
\sqrt{2} = \frac{p}{q}.
\]
Moreover, we may choose such integers $p$ and $q$ so that the fraction $\frac{p}{q}$ is in lowest terms, meaning that
\[
\gcd(p,q) = 1.
\]

Squaring both sides of
\[
\sqrt{2} = \frac{p}{q}
\]
gives
\[
2 = \frac{p^2}{q^2}.
\]
Multiplying both sides by $q^2$ yields
\[
2q^2 = p^2.
\]
Therefore $p^2$ is even.

We now show that if an integer has even square, then the integer itself is even. Suppose $n$ is an integer and $n$ is odd. Then there exists an integer $k$ such that
\[
n = 2k+1.
\]
Hence
\[
n^2 = (2k+1)^2 = 4k^2+4k+1 = 2(2k^2+2k)+1,
\]
which is odd. Thus an odd integer has odd square. Consequently, if $n^2$ is even, then $n$ must be even.

Since $p^2$ is even, it follows that $p$ is even. Hence there exists an integer $r$ such that
\[
p = 2r.
\]
Substituting this into the equation $2q^2 = p^2$, we obtain
\[
2q^2 = (2r)^2 = 4r^2.
\]
Dividing both sides by $2$ gives
\[
q^2 = 2r^2.
\]
Therefore $q^2$ is even. By the same argument as above, $q$ must be even.

Thus both $p$ and $q$ are even, so $2$ divides both $p$ and $q$. Therefore
\[
\gcd(p,q) \geq 2,
\]
which contradicts the assumption that
\[
\gcd(p,q) = 1.
\]

This contradiction shows that $\sqrt{2}$ cannot be rational. Therefore $\sqrt{2}$ is irrational.
\end{proof}
```